\documentclass[10pt,a4paper]{amsart}
\usepackage{amsmath,amssymb,amsthm,graphicx}
% Retained minimal preamble for the Stacks Project chapter more-algebra.tex (More on Algebra).
% The chapter file pinned-original-source.tex carries only \input{preamble} on its first line, so no
% preamble text lives inside the chapter itself.  The Stacks Project uses one single shared preamble for
% every chapter; the snapshot stages it identically as sources/stacks-chow/preamble.tex and
% sources/stacks-etale-cohomology/preamble.tex.  The block below is copied mechanically, line for line,
% from sources/stacks-chow/preamble.tex lines 14-18 (the Xy-pic setup) and 154-217 (the native theorem
% declarations together with the mathematical macros).  Nothing was added, removed or redefined; the
% wrapper class and the amsmath/amssymb/amsthm/graphicx base packages are supplied by the assembler.
% Chapter: more-algebra.tex -- More on Algebra.
% For commutative diagrams we use Xy-pic
\usepackage[all]{xy}

% We use 2cell for 2-commutative diagrams.
\xyoption{2cell}

% Theorem environments.
%
\theoremstyle{plain}
\newtheorem{theorem}[subsection]{Theorem}
\newtheorem{proposition}[subsection]{Proposition}
\newtheorem{lemma}[subsection]{Lemma}

\theoremstyle{definition}
\newtheorem{definition}[subsection]{Definition}
\newtheorem{example}[subsection]{Example}
\newtheorem{exercise}[subsection]{Exercise}
\newtheorem{situation}[subsection]{Situation}

\theoremstyle{remark}
\newtheorem{remark}[subsection]{Remark}
\newtheorem{remarks}[subsection]{Remarks}

\numberwithin{equation}{subsection}

% Macros
%
\def\lim{\mathop{\mathrm{lim}}\nolimits}
\def\colim{\mathop{\mathrm{colim}}\nolimits}
\def\Spec{\mathop{\mathrm{Spec}}}
\def\Hom{\mathop{\mathrm{Hom}}\nolimits}
\def\Ext{\mathop{\mathrm{Ext}}\nolimits}
\def\SheafHom{\mathop{\mathcal{H}\!\mathit{om}}\nolimits}
\def\SheafExt{\mathop{\mathcal{E}\!\mathit{xt}}\nolimits}
\def\Sch{\mathit{Sch}}
\def\Mor{\mathop{\mathrm{Mor}}\nolimits}
\def\Ob{\mathop{\mathrm{Ob}}\nolimits}
\def\Sh{\mathop{\mathit{Sh}}\nolimits}
\def\NL{\mathop{N\!L}\nolimits}
\def\CH{\mathop{\mathrm{CH}}\nolimits}
\def\proetale{{pro\text{-}\acute{e}tale}}
\def\etale{{\acute{e}tale}}
\def\QCoh{\mathit{QCoh}}
\def\Ker{\mathop{\mathrm{Ker}}}
\def\Im{\mathop{\mathrm{Im}}}
\def\Coker{\mathop{\mathrm{Coker}}}
\def\Coim{\mathop{\mathrm{Coim}}}

% Boxtimes
%
\DeclareMathSymbol{\boxtimes}{\mathbin}{AMSa}{"02}

%
% Macros for moduli stacks/spaces
%
\def\QCohstack{\mathcal{QC}\!\mathit{oh}}
\def\Cohstack{\mathcal{C}\!\mathit{oh}}
\def\Spacesstack{\mathcal{S}\!\mathit{paces}}
\def\Quotfunctor{\mathrm{Quot}}
\def\Hilbfunctor{\mathrm{Hilb}}
\def\Curvesstack{\mathcal{C}\!\mathit{urves}}
\def\Polarizedstack{\mathcal{P}\!\mathit{olarized}}
\def\Complexesstack{\mathcal{C}\!\mathit{omplexes}}
% \Pic is the operator that assigns to X its picard group, usage \Pic(X)
% \Picardstack_{X/B} denotes the Picard stack of X over B
% \Picardfunctor_{X/B} denotes the Picard functor of X over B
\def\Pic{\mathop{\mathrm{Pic}}\nolimits}
\def\Picardstack{\mathcal{P}\!\mathit{ic}}
\def\Picardfunctor{\mathrm{Pic}}
\def\Deformationcategory{\mathcal{D}\!\mathit{ef}}
\usepackage{cleveref}
\crefname{originallemma}{original lemma}{original lemmas}
\Crefname{originallemma}{Original lemma}{Original lemmas}
\makeatletter
\newlabel{more-algebra-lemma-principal-fitting-ideal}{{080Z}{source}{}{sourcepointer-1}{}}
\newlabel{more-algebra-lemma-principal-fitting-ideal@cref}{{[originallemma][0][]080Z}{[1][0][]source}{}{}{}}
\makeatother
\crefname{originallemma}{original lemma}{original lemmas}
\Crefname{originallemma}{Original lemma}{Original lemmas}
\makeatletter
\newlabel{algebra-lemma-NAK}{{00DV}{source}{}{sourcepointer-2}{}}
\newlabel{algebra-lemma-NAK@cref}{{[originallemma][0][]00DV}{[1][0][]source}{}{}{}}
\makeatother
\crefname{originallemma}{original lemma}{original lemmas}
\Crefname{originallemma}{Original lemma}{Original lemmas}
\makeatletter
\newlabel{algebra-lemma-matrix-right-inverse}{{080R}{source}{}{sourcepointer-3}{}}
\newlabel{algebra-lemma-matrix-right-inverse@cref}{{[originallemma][0][]080R}{[1][0][]source}{}{}{}}
\makeatother
\title{A finite module whose k-th Fitting ideal is principal is generated by k elements after removing its f-torsion}
\author{The Stacks Project Authors}
\date{stacks-project:more-algebra.tex:local-snapshot-20261012; GNU FDL 1.2 or later}
\begin{document}
\maketitle
\setcounter{subsection}{1}
\noindent\textit{Editorial note.} All delivered mathematical text is the author's own, copied from the pinned Stacks Project chapter file more-algebra.tex; no statement or proof step was rewritten, reordered or removed.\par\medskip
\noindent\textit{Editorial note.} The chapter file pinned-original-source.tex carries only \textbackslash{}input\{preamble\} on its first line, so no preamble text lives inside the chapter. The retained file sources/stacks-more-algebra/retained-preamble.tex is a mechanical, line-for-line copy of the project's single shared preamble as staged in the snapshot (Xy-pic setup plus theorem declarations and mathematical macros). Nothing was added, removed or redefined, and the excerpt serves only as preamble, never as a span coordinate.\par\medskip
\noindent\textit{Editorial note.} Both counted spans start at the \textbackslash{}begin line of their environment, so the statement and the author's proof are delivered whole; surrounding prose of the section is deliberately outside the counted spans and is carried as standard prerequisites instead.\par\medskip
\noindent\textit{Editorial note.} The citation fields (source\_id, source\_version, authors, source\_url, license.name, source\_title) are identical to those of every other card this lane has produced from this chapter, so the batch selector treats them as one source.\par\medskip
\noindent\textit{Allowed standard prerequisites (editorial).}\begin{itemize}
\item Fitting ideal of a finite module and its description through the minors of a presentation matrix. defined in the same section, source lines 1320-1578, before the counted statement. Supplies the ideal whose being principal is the hypothesis.
\item Nakayama's lemma in the form that a set of elements generating a principal ideal over a local ring contains a generator. the author cites Algebra, Lemma algebra-lemma-NAK at proof lines 1591-1593. Allows the Fitting ideal to be generated by a single chosen minor.
\item Existence of a right inverse for a matrix with an invertible minor, up to the factor given by that minor. the author cites Algebra, Lemma algebra-lemma-matrix-right-inverse at proof lines 1601-1603. Produces the matrix whose columns give the combinations annihilated by f.
\end{itemize}\medskip
\ifdefined\hypertarget\hypertarget{sourcepointer-1}{}\fi
\noindent\textit{Source reference.} Original lemma 080Z: focused-seeds/AGP/sources/stacks-snapshot-20261011/tags.txt, source lines 10199--10199. The lemma's own canonical tag, quoted so that readers can resolve the label inside the delivered text; it points at the counted statement itself and carries no proof content.\par\medskip
\ifdefined\hypertarget\hypertarget{sourcepointer-2}{}\fi
\noindent\textit{Source reference.} Original lemma 00DV: focused-seeds/AGP/sources/stacks-snapshot-20261011/tags.txt, source lines 512--512. Navigation pointer to a statement the counted proof refers to but does not contain; nothing load-bearing is delegated to it.\par\medskip
\ifdefined\hypertarget\hypertarget{sourcepointer-3}{}\fi
\noindent\textit{Source reference.} Original lemma 080R: focused-seeds/AGP/sources/stacks-snapshot-20261011/tags.txt, source lines 10191--10191. Navigation pointer to a statement the counted proof refers to but does not contain; nothing load-bearing is delegated to it.\par\medskip
\begin{lemma}
\label{lemma-principal-fitting-ideal}
Let $R$ be a local ring. Let $M$ be a finite $R$-module. Let $k \geq 0$.
Assume that $\text{Fit}_k(M) = (f)$ for some $f \in R$.
Let $M'$ be the quotient of $M$ by $\{x \in M \mid fx = 0\}$. Then
$M'$ can be generated by $k$ elements.
\end{lemma}

\begin{proof}
Choose generators $x_1, \ldots, x_n \in M$ corresponding to the
surjection $R^{\oplus n} \to M$. Since $R$ is local if a set
of elements $E \subset (f)$ generates $(f)$, then some $e \in E$ generates
$(f)$, see Algebra, Lemma \ref{algebra-lemma-NAK}. Hence we may pick
$z_1, \ldots, z_{n - k}$ in the kernel of $R^{\oplus n} \to M$ such
that some $(n - k) \times (n - k)$ minor of the $n \times (n - k)$
matrix $A = (z_{ij})$ generates $(f)$. After renumbering the $x_i$ we may
assume the first minor $\det(z_{ij})_{1 \leq i, j \leq n - k}$
generates $(f)$, i.e., $\det(z_{ij})_{1 \leq i, j \leq n - k} = uf$
for some unit $u \in R$. Every other minor is a multiple of $f$.
By Algebra, Lemma \ref{algebra-lemma-matrix-right-inverse} there exists a
$n - k \times n - k$ matrix $B$ such that
$$
AB = f
\left(
\begin{matrix}
u 1_{n - k \times n - k} \\
C
\end{matrix}
\right)
$$
for some matrix $C$ with coefficients in $R$. This implies that for every
$i \leq n - k$ the element $y_i = ux_i + \sum_j c_{ji}x_j$ is annihilated
by $f$. Since $M/\sum Ry_i$ is generated by the images of
$x_{n - k + 1}, \ldots, x_n$ we win.
\end{proof}

\end{document}
